%----------------------------------------------------------------------------------------
%	PACKAGES AND OTHER DOCUMENT CONFIGURATIONS
%----------------------------------------------------------------------------------------

\documentclass[11pt,a4paper,oneside]{memoir} % Change font size here (allowable values are 9pt-12pt), change the paper size, specify one or two sided printing and specify whether to show trimming lines

\input{structure.tex} % Include the file containing the code defining the structure and style of the document

%------------------------------------------------
% Thesis Information
\newcommand{\themonth}{September 2025} % The month of writing    

\title{Development of an Optical Tracking System for Real-time Robotic Applications} % Thesis title

\author{Shannon Pitman} 

\date{}

% Supervisor
\newcommand{\supervisor}{Arnold Pretorius}

\newcommand{\institution}{University of Cape Town\xspace} % University/institution name

\newcommand{\researchgroup}{\href{http://mechatronicsystems.group}{MechatronicsSystems.Group}\xspace} % University/institution name

\newcommand{\department}{Department of Mechanical Engineering\xspace} 

\newcommand{\degreetitle}{BSc (Eng) in Mechanical and Mechatronic Engineering}

%------------------------------------------------
% % Fonts
% \renewcommand*{\acffont}[1]{{\normalsize\itshape #1}} % Font style for the acronym text (e.g. Do It Yourself)
% \renewcommand*{\acfsfont}[1]{{\normalsize\upshape #1}} % Font style for the acronym in bracket (e.g. (DIY))

%------------------------------------------------
% Hyphenations

\hyphenation{a-no-ma-lous a-no-ma-ly amounts breaches} % Specify custom hyphenation points in words with dashes where you would like hyphenation to occur, or alternatively, don't put any dashes in a word to stop hyphenation altogether

%----------------------------------------------------------------------------------------
%	TITLE PAGE
%----------------------------------------------------------------------------------------

\renewcommand{\maketitlehooka}{}

\renewcommand{\maketitlehookb}{
\centering
\vfill
\includegraphics[width=0.6\linewidth]{Figures/UCT_MSG_wide.png}\\[0.5cm] % Institution logo
\par
\hrulefill}


\renewcommand{\maketitlehookc}{
% People details left, institution info right
\vspace{0.5cm}
\noindent
\begin{tabular*}{\textwidth}{@{} l @{\extracolsep{\fill}} r}
    \textbf{Author:} \author &  \researchgroup \\
    \textbf{Supervisor:} \supervisor & \department\\
    \textbf{Date:} \themonth & \institution\\
\end{tabular*}

\vspace{1cm}

\emph{Submitted in partial fulfilment of the requirements for the degree of} \\
\degree
}

%----------------------------------------------------------------------------------------

\makeindex % Write an index file

\begin{document}

\begin{titlingpage}
\maketitle % Print the title page
\end{titlingpage}

\frontmatter % Use roman page numbering style (i, ii, iii, iv...) for the pre-content pages

%----------------------------------------------------------------------------------------
%	ABSTRACT
%----------------------------------------------------------------------------------------

\chapterstyle{default} % Apply thesis chapter style
\chapter*{Abstract}
%\addcontentsline{toc}{chapter}{Abstract} % Add to table of contents

Engineering students often need a clear, practical model for writing dissertations; this document provides an overview tailored to final‑year mechatronics projects and demonstrates effective use of a LaTeX thesis template. Many student reports lack consistent structure, rigorous scoping, and appropriate use of figures, tables, code listings, and citations, which impedes readability and assessment.

This work contributes an exemplar that explains what each report section should contain, offers typical subheadings, and embeds working examples of template elements to copy and adapt. The document is implemented in XeLaTeX using the \texttt{memoir} class with chapters for Introduction, Literature Review, Methodology, Results, Discussion, Conclusions, Recommendations, and Appendices, along with acronyms and bibliography support; the result is a shareable guide that doubles as a starting point for students to produce professional, evidence‑based engineering reports.
\vspace*{1cm}

\subsection*{\emph{How to write an abstract in 5 minutes?}}

\begin{flushright}
\textit{Answer the following questions in narrative form with no more than two sentences per question:}

\begin{itemize}
    \item \textit{What is the background to the project?}
    \item \textit{What is the problem that is being solved?}
    \item \textit{What is the contribution of the project to this problem?}
    \item \textit{What is the main conclusion of the project?}
\end{itemize}
\end{flushright}


\cleartoverso % Force a break to an even page

%----------------------------------------------------------------------------------------
%	TABLE OF CONTENTS
%----------------------------------------------------------------------------------------

%\renewcommand{\contentsname}{\chapnamefont\fontsize{14}{14}\selectfont Contents}
\tableofcontents* % Print the table of contents

\cleartoverso % Force a break to an even page

%----------------------------------------------------------------------------------------
%	LIST OF FIGURES
%----------------------------------------------------------------------------------------

%\renewcommand{\listfigurename}{\chapnamefont\fontsize{14}{14}\selectfont List of Figures}
\listoffigures % Print the list of figures

\cleartoverso % Force a break to an even page

%----------------------------------------------------------------------------------------
%	LIST OF TABLES
%----------------------------------------------------------------------------------------

%\renewcommand{\listtablename}{\chapnamefont\fontsize{14}{14}\selectfont List of Tables}
\listoftables % Print the list of tables

\cleartoverso % Force a break to an even page

%----------------------------------------------------------------------------------------
%	LIST OF EQUATIONS
%----------------------------------------------------------------------------------------

% Create a new list for equations (memoir syntax)
\newlistof{listofequations}{equ}{List of Equations}
\newlistentry{myequation}{equ}{0}
\cftsetindents{myequation}{0em}{2.5em}

\newcommand{\myequation}[1]{%
  \refstepcounter{myequation}
  \addcontentsline{equ}{myequation}{\protect\numberline{\theequation}#1}}
\cleartoverso % Force a break to an even page

%----------------------------------------------------------------------------------------
%	LIST OF ACRONYMS
%----------------------------------------------------------------------------------------

\printglossary[type=abbreviations]
\printglossary[type=symbols]
\printglossary[type=main]
\cleartoverso % Force a break to an even page
%----------------------------------------------------------------------------------------
%	CONTENT CHAPTERS
%----------------------------------------------------------------------------------------

\mainmatter % Begin numeric (1,2,3...) page numbering

\chapterstyle{thesis} % Change the style of the Chapter header to that defined in structure.tex

\pagestyle{Ruled} % Include the chapter/section in the header along with a horizontal rule underneath

\include{Chapters/Introduction} % Include the introduction chapter
\include{Chapters/Literature_Review} % Literature Review of overarching goal
\chapter{System Specifications}
\label{Systems_Specifications}

\section{Introduction}

This chapter translates the scope statements and design objectives from Chapter~\ref{Introduction} into measurable, verifiable specifications for two subsystems: the \gls{OMC} subsystem (OpenMV RT1062 camera network and supporting software) and the \gls{OCP} subsystem (\gls{GA} and cost function). Together they form the end-to-end pipeline: the \gls{OCP} algorithm determines the optimal camera configuration offline, and the resulting arrangement is deployed for real-time marker tracking by the \gls{OMC} subsystem.

Following NASA systems engineering practice \parencite{nasaSE2017}, each requirement carries a unique identifier, an importance level, and two acceptance thresholds: a \textit{marginal} value (minimum acceptable performance) and an \textit{ideal} value (target under well-optimised conditions). Importance levels are defined as follows:

\begin{itemize}
    \item \textbf{Critical (C)}: The system cannot fulfil its primary function without meeting this requirement.
    \item \textbf{High (H)}: Failure to meet this requirement significantly degrades system effectiveness.
    \item \textbf{Medium (M)}: Desirable for optimal performance; the system remains operational at reduced capacity without it.
    \item \textbf{Low (L)}: Beneficial but not essential; the system meets its primary objectives regardless.
\end{itemize}

Each requirement carries a traceability code: \texttt{S-\textit{XX}} codes refer to scope specifications from Chapter~\ref{Introduction}; \texttt{R-\textit{XX}} codes identify external sources (datasheets, validated literature results, or engineering conventions) that ground the stated thresholds. Requirements are divided into two categories. \textit{Performance} requirements define quantitative thresholds that the system must meet during operation. \textit{Functional} requirements describe the operational capabilities, constraints, and interface properties the system must possess; these are verified through demonstration or inspection rather than continuous measurement.The full mapping is given in Table~\ref{tab:trace-refs}.

Table~\ref{tab:metrics-overview} provides a consolidated overview of all specified metrics; details are provided in the subsystem sections below.

{\small\setlength{\tabcolsep}{4pt}
\begin{longtable}{>{\raggedright\arraybackslash}p{1.2cm} >{\raggedright\arraybackslash}p{2.6cm} >{\raggedright\arraybackslash}p{1.1cm} >{\raggedright\arraybackslash}p{0.9cm} >{\raggedright\arraybackslash}p{\dimexpr\linewidth-1.2cm-2.6cm-1.1cm-0.9cm-10\tabcolsep\relax}}
\caption[Overview of all system specification metrics.]{Overview of all system specification metrics. Category key: P\,=\,Performance, F\,=\,Functional.}
\label{tab:metrics-overview}\\
\toprule
\textbf{ID} & \textbf{Metric} & \textbf{Sub-system} & \textbf{Cat.} & \textbf{Purpose} \\
\midrule
\endfirsthead
\multicolumn{5}{l}{\small\textit{(Continued from previous page)}}\\[2pt]
\toprule
\textbf{ID} & \textbf{Metric} & \textbf{Sub-system} & \textbf{Cat.} & \textbf{Purpose} \\
\midrule
\endhead
\midrule
\multicolumn{5}{r}{\small\textit{(Continued on next page)}}\\
\endfoot
\bottomrule
\endlastfoot
OMC-P-01 & Sensor Resolution        & OMC & P & Spatial sampling capability of each camera; governs minimum detectable marker size in the image plane. Higher resolution improves localisation but increases latency and reduces frame rate. \\[4pt]
OMC-P-02 & 3D Position Accuracy     & OMC & P & Quantifies positional error against a known ground-truth reference (OptiTrack); sets the foundation for trajectory estimate quality. \\[4pt]
OMC-P-03 & Frame Rate               & OMC & P & Frequency of image capture must be sufficient to sample fast robotic manoeuvres without aliasing. \\[4pt]
OMC-P-04 & System Latency           & OMC & P & Pipeline delay from image capture to pose output; critical for closed-loop \gls{EKF} that perform optimally with low-latency measurements. \\[4pt]
OMC-P-05 & Calibration Residual     & OMC & P & Mean ray-convergence offset reported by EasyWand6; a direct indicator of triangulation precision for a given arrangement. \\[4pt]
OMC-P-06 & Marker Reconstruction Completeness & OMC & P & Percentage of captured frames in which all markers are successfully reconstructed; reflects network coverage and redundancy. \\
\midrule
OMC-F-01 & Capture Volume           & OMC & F & Minimum observable 3D volume within which all performance requirements are simultaneously satisfied. \\[4pt]
OMC-F-02 & Maximum Detection Range  & OMC & F & Furthest camera-to-marker distance at which a marker can be reliably detected given lens type and marker size. \\[4pt]
OMC-F-03 & Camera Count             & OMC & F & Minimum number of cameras required for multi-view redundancy, 3D reconstruction, and budget. \\[4pt]
OMC-F-04 & Ease of Setup            & OMC & F & Constrains the system to fit within the available laboratory space without structural modification. \\[4pt]
OMC-F-05 & System Cost              & OMC & F & Restricts total hardware and software expenditure to within the available research budget. \\[4pt]
OMC-F-06 & Portability              & OMC & F & Ease of disassembly and redeployment without specialist equipment. \\[4pt]
OMC-F-07 & Battery Life             & OMC & F & Duration of continuous tracking available before recharge when operating with wireless power transmission. \\
\midrule
OCP-P-01 & Workspace Coverage       & OCP & P & Fraction of target grid points observed by $\geq 2$ cameras; necessary precondition for triangulation. \\[4pt]
OCP-P-02 & Triangulation Angle Compliance & OCP & P & Fraction of target points where $\geq 1$ camera pair satisfies $40\degree \leq \theta \leq 140\degree$ for well-conditioned reconstruction. \\[4pt]
OCP-P-03 & Combined GA Objective Cost & OCP & P & Scalar $J$ (Eq.~\ref{eq:combined-cost}) directly minimised by the \gls{GA}; combines the normalised resolution uncertainty cost $J_{\text{res}}$ and occlusion robustness cost $J_{\text{occ}}$ into a single dimensionless measure of placement quality. Values below 1.0 indicate improvement over a random baseline. \\[4pt]
OCP-P-04 & Optimisation Runtime     & OCP & P & Overall time for the GA to complete a full run; bounds the pre-deployment requisite duration. \\[4pt]
OCP-P-05 & Calibration Residual Improvement & OCP & P & Percentage reduction in the OpenMV system calibration residual vs.\ the ad-hoc baseline, with OptiTrack as ground-truth reference; primary experimental validation metric. \\[4pt]
OCP-P-06 & Marker Reconstruction Improvement & OCP & P & Percentage-point increase in reconstruction completeness vs.\ ad-hoc baseline; validates real-world tracking gains. \\
\midrule
OCP-F-01 & Solution Repeatability   & OCP & F & Consistency of the optimised cost across independent runs; characterises stochastic variance. \\[4pt]
OCP-F-02 & Output Interface         & OCP & F & Format in which the optimal configuration is exported for physical implementation. \\[4pt]
OCP-F-03 & Objective Weight Parameterizability & OCP & F & User-adjustable weighting of resolution and occlusion objectives in the combined cost function. \\
\end{longtable}}

{\small\setlength{\tabcolsep}{4pt}
\begin{longtable}{p{1.0cm} p{1.5cm} >{\raggedright\arraybackslash}p{\dimexpr\linewidth-1.0cm-1.5cm-6\tabcolsep\relax}}
\caption{Reference traceability codes used in the subsystem specification tables.}
\label{tab:trace-refs}\\
\toprule
\textbf{Code} & \textbf{Type} & \textbf{Source} \\
\midrule
\endfirsthead
\toprule
\textbf{Code} & \textbf{Type} & \textbf{Source} \\
\midrule
\endhead
\bottomrule
\endlastfoot
S-01 & Internal & \textit{Ch.\ref{Introduction}}: The system shall provide real-time 3D pose tracking of robotic markers within an indoor laboratory testing arena. \\[4pt]
S-02 & Internal & \textit{Ch.\ref{Introduction}}: The system shall achieve tracking accuracy and reconstruction completeness sufficient to support closed-loop \gls{UAV} flight control. \\[4pt]
S-03 & Internal & \textit{Ch.\ref{Introduction}}: The \gls{OCP} subsystem shall produce a camera configuration that yields measurably better tracking performance than an ad-hoc arrangement. \\[4pt]
S-04 & Internal & \textit{Ch.\ref{Introduction}}: The system shall be practically deployable within the constraints of the MechatronicSystems.Group laboratory (space, budget, available persons). \\[4pt]
R-01 & External & OptiTrack camera specification sheet \parencite{OptiTrackDatasheet}. \\[4pt]
R-02 & External & Aurand et al.\ (2017): accuracy mapping of a 42-camera OptiTrack Prime 41 system; 97\,\% of capture area below 0.2\,mm error \parencite{aurandAccuracyMapOptical2017}. \\[4pt]
R-03 & External & Nagymát{\'e} \& Kiss (2018): review of OptiTrack validation studies; identifies calibration residual thresholds of $<0.5$\,mm (acceptable) and $<0.3$\,mm (excellent) \parencite{nagymateMotionCaptureSystem2018}. \\[4pt]
R-04 & External & Rahimian \& Kearney (2017): physical OptiTrack deployment achieved 86\,\% marker reconstruction (optimised) vs.\ 30\,\% (ad-hoc) \parencite{rahimianOptimalCameraPlacement2017}. \\[4pt]
R-05 & External & Ziegler et al.\ (2023): triangulation error grows exponentially as convergence angle approaches $0\degree$ or $180\degree$; well-conditioned range $40\degree$--$140\degree$ \parencite{zieglerOptimalCameraPlacement2023}. \\
\end{longtable}}

\section{Overall System Architecture}

The system comprises two coupled subsystems, illustrated in Figure \ref{fig:system-block-diagram}. The \gls{OCP} subsystem runs \textit{offline}: it accepts workspace geometry and camera hardware specifications and outputs an optimal camera configuration via the genetic algorithm of Chapter \ref{OCP}, which is then physically deployed. The \gls{OMC} subsystem operates \textit{online}: the OpenMV RT1062 camera network captures infrared images of passive retroreflective markers on the tracked platform; triangulation software reconstructs 3D marker positions; and the resulting pose estimates are either streamed to a robotic controller or collated for data analysis. The two subsystems are coupled through the camera configuration---the geometric arrangement that the \gls{OCP} algorithm optimises.

\begin{figure}[htbp]
\centering
\begin{tikzpicture}[
    node distance=1.1cm and 2.0cm,
    process/.style={rectangle, rounded corners=3pt, minimum width=3.4cm, minimum height=0.75cm, align=center, draw=black!70, fill=blue!8, font=\small},
    hardware/.style={rectangle, minimum width=3.4cm, minimum height=0.75cm, align=center, draw=black!70, fill=orange!15, font=\small},
    datanode/.style={rectangle, minimum width=3.0cm, minimum height=0.65cm, align=center, draw=black!70, fill=green!12, font=\small},
    externalnode/.style={rectangle, rounded corners=3pt, minimum width=3.4cm, minimum height=0.75cm, align=center, draw=black!70, fill=gray!12, font=\small},
    arrow/.style={-{Stealth[length=1.5mm]}, thick},
    dashedarrow/.style={-{Stealth[length=2.5mm]}, thick, dashed, gray!70},
    grouplabel/.style={font=\small\bfseries, text=black!70},
]
\node (specs)  [process]                                        {Workspace \& Camera\\Specifications};
\node (ga)     [process, below=of specs]                        {GA Optimiser\\(Ch. \ref{OCP})};
\node (config) [datanode, below=of ga]                          {Optimal Camera\\Configuration};
\draw[arrow] (specs)  -- (ga);
\draw[arrow] (ga)     -- (config);

\node (mount)  [hardware, right=4.5cm of specs] {7$\times$ OpenMV RT1062 Cameras\\(IR illuminators + sensors)};
\node (software) [process,  below=of mount]                       {EKF Software\\(Calibration \& 3D Reconstruction)};
\node (pose)   [datanode, below=of software]                      {3D Pose Estimate};
\node (ctrl)   [externalnode, below=of pose]                    {Data Consumer\\(Controller / Logger)};
\draw[arrow] (mount)  -- (software);
\draw[arrow] (software) -- (pose);
\draw[arrow] (pose)   -- (ctrl);

\draw[dashedarrow] (config.east) -- node[above, font=\scriptsize, text=black] {Physical deployment} (mount.west);

\node[grouplabel, above=0.5cm of specs] {OCP Subsystem (Offline)};
\node[grouplabel, above=0.5cm of mount] {OMC Subsystem (Real-Time)};

\begin{scope}[on background layer]
    \node[draw=blue!30, fill=blue!4, rounded corners=5pt,
          fit=(specs)(ga)(config), inner sep=0.45cm] {};
    \node[draw=orange!40, fill=orange!4, rounded corners=5pt,
          fit=(mount)(software)(pose)(ctrl), inner sep=0.45cm] {};
\end{scope}
\end{tikzpicture}
\caption[Overall system architecture block diagram.]{Overall system architecture. The \gls{OCP} subsystem (left) determines the optimal camera placement, which the informs the physical arrangement of the cameras. The \gls{OMC} subsystem (right) operates in real time to stream or log 3D pose estimates for use by a controller or for later data analysis.}
\label{fig:system-block-diagram}
\end{figure}


\section{OMC Subsystem Requirements}

The \gls{OMC} subsystem captures and reconstructs 3D marker positions within the tracking arena using OpenMV RT1062 cameras. The system is robot-agnostic: it tracks any platform carrying passive retroreflective markers, with \gls{UAV}s as the primary test case. Performance is validated against an OptiTrack system running Motive, which provides the ground-truth reference. Specifications are informed by the demands of closed-loop robotic control (low latency, high frame rate, positional accuracy) \parencite{aurandAccuracyMapOptical2017, nagymateMotionCaptureSystem2018} and by the OpenMV RT1062 hardware characteristics \parencite{OpenMVCamRT1062a}.

\subsection{Performance Requirements}

Table \ref{tab:omc-perf-desc} defines each metric; the quantitative thresholds are given in Table \ref{tab:omc-perf}.

{\small\setlength{\tabcolsep}{4pt}
\begin{longtable}{p{1.5cm} >{\raggedright\arraybackslash}p{3.0cm} >{\raggedright\arraybackslash}p{\dimexpr\linewidth-1.5cm-3.0cm-6\tabcolsep\relax}}
\caption[OMC performance metric definitions.]{OMC performance metric definitions. Quantitative thresholds are given in Table \ref{tab:omc-perf}.}
\label{tab:omc-perf-desc}\\
\toprule
\textbf{ID} & \textbf{Metric} & \textbf{Description} \\
\midrule
\endfirsthead
\multicolumn{3}{l}{\small\textit{(Continued from previous page)}}\\[2pt]
\toprule
\textbf{ID} & \textbf{Metric} & \textbf{Description} \\
\midrule
\endhead
\bottomrule
\endlastfoot
OMC-P-01 & Sensor Resolution & Minimum imaging array resolution per camera, governing the detection capabilities of markers in the image plane. \\[4pt]
OMC-P-02 & 3D Position Accuracy & \gls{RMS} error of reconstructed marker positions against OptiTrack ground-truth measurements across the capture volume. \\[4pt]
OMC-P-03 & Frame Rate & Sampling rate of the camera system must be sufficient to capture maximum tracked-platform velocities without spatial aliasing or motion blur. \\[4pt]
OMC-P-04 & System Latency & End-to-end pipeline delay from image acquisition to 3D pose availability at the downstream system (controller or data logger). Directly constrains the bandwidth of any feedback loop using \gls{OMC} state estimates. \\[4pt]
OMC-P-05 & Calibration Residual & Mean \gls{RMS} offset (mm) reported by EasyWand6 following the wand calibration procedure. Reflects geometric accuracy of the camera network and triangulation precision. \\[4pt]
OMC-P-06 & Marker Reconstruction Completeness & Percentage of captured frames in which all markers in the rigid body are successfully resolved. Values below the marginal threshold render continuous trajectory estimation unreliable. \\
\end{longtable}}

{\small\setlength{\tabcolsep}{4pt}
\begin{longtable}{p{2.5cm} p{1.8cm} p{0.9cm} p{2.0cm} >{\raggedright\arraybackslash}p{\dimexpr(\linewidth-1.5cm-1.8cm-0.7cm-1.0cm-12\tabcolsep)/2\relax} >{\raggedright\arraybackslash}p{\dimexpr(\linewidth-1.5cm-1.8cm-0.7cm-1.0cm-12\tabcolsep)/2\relax}}
\caption[OMC subsystem performance specifications.]{OMC subsystem performance specifications. Importance: C\,=\,Critical, H\,=\,High, M\,=\,Medium, L\,=\,Low. Metric definitions are given in Table~\ref{tab:omc-perf-desc}.}
\label{tab:omc-perf}\\
\toprule
\textbf{ID} & \textbf{Trace} & \textbf{Imp.} & \textbf{Units} & \textbf{Marginal} & \textbf{Ideal} \\
\midrule
\endfirsthead
\multicolumn{6}{l}{\small\textit{(Continued from previous page)}}\\[2pt]
\toprule
\textbf{ID} & \textbf{Trace} & \textbf{Imp.} & \textbf{Units} & \textbf{Marginal} & \textbf{Ideal} \\
\midrule
\endhead
\midrule
\multicolumn{6}{r}{\small\textit{(Continued on next page)}}\\
\endfoot
\bottomrule
\endlastfoot
OMC-P-01 & S-01, R-01 & H & Pixels & $\geq QQVGA (160 \times 120)$ & $ VGA (640 \times 480)$ \\[4pt]
OMC-P-02 & S-02, R-02 & C & cm\,(RMS) & $< 1.0$ & $< 0.3$ \\[4pt]
OMC-P-03 & S-02, R-01 & H & fps & $\geq 40$ & $\geq 60$ \\[4pt]
OMC-P-04 & S-02, R-01 & H & ms & $< 10$ & $< 5.5$ \\[4pt]
OMC-P-05 & S-02, R-03 & H & mm & $< 0.7$ & $< 0.3$ \\[4pt]
OMC-P-06 & S-02, R-04 & C & \% & $> 85$ & $> 99$ \\
\end{longtable}}

\subsection{Functional Requirements}

Table~\ref{tab:omc-func-desc} defines each metric; the quantitative thresholds are given in Table~\ref{tab:omc-func}.

{\small\setlength{\tabcolsep}{4pt}
\begin{longtable}{p{1.5cm} >{\raggedright\arraybackslash}p{3.0cm} >{\raggedright\arraybackslash}p{\dimexpr\linewidth-1.5cm-3.0cm-6\tabcolsep\relax}}
\caption[OMC functional metric definitions.]{OMC functional metric definitions. Quantitative thresholds are given in Table~\ref{tab:omc-func}.}
\label{tab:omc-func-desc}\\
\toprule
\textbf{ID} & \textbf{Metric} & \textbf{Description} \\
\midrule
\endfirsthead
\multicolumn{3}{l}{\small\textit{(Continued from previous page)}}\\[2pt]
\toprule
\textbf{ID} & \textbf{Metric} & \textbf{Description} \\
\midrule
\endhead
\bottomrule
\endlastfoot
OMC-F-01 & Capture Volume & Usable 3D volume in which all performance requirements (OMC-P-01 to OMC-P-06) are simultaneously satisfied. Defines the tracking volume available to the robotic platform. \\[4pt]
OMC-F-02 & Maximum Detection Range & Maximum camera-to-marker distance at which a standard 15.9\,mm retroreflective marker can be reliably detected for each lens type. \\[4pt]
OMC-F-03 & Camera Count & Minimum number of cameras required to simultaneously achieve the prescribed coverage and accuracy requirements across the full capture volume. \\[4pt]
OMC-F-04 & Setup Space & Maximum floor area and ceiling height required for system installation. Must not exceed the dimensions of the available laboratory space. \\[4pt]
OMC-F-05 & System Cost & Total expenditure on camera hardware, cabling, networking, and licensing. Must remain within the available research project budget and be low-cost in comparison to proprietary systems. \\[4pt]
OMC-F-06 & Portability & Ease with which the camera network can be disassembled and redeployed. \\[4pt]
OMC-F-07 & Battery Life & Duration of continuous tracking available before recharge when the system is operating with wireless power transmission rather than wired powered supply. \\
\end{longtable}}

{\small\setlength{\tabcolsep}{4pt}
\begin{longtable}{p{1.5cm} p{1.8cm} p{0.7cm} p{1.0cm} >{\raggedright\arraybackslash}p{\dimexpr(\linewidth-1.5cm-1.8cm-0.7cm-1.0cm-12\tabcolsep)/2\relax} >{\raggedright\arraybackslash}p{\dimexpr(\linewidth-1.5cm-1.8cm-0.7cm-1.0cm-12\tabcolsep)/2\relax}}
\caption[OMC subsystem functional specifications.]{OMC subsystem functional specifications. Metric definitions are given in Table~\ref{tab:omc-func-desc}.}
\label{tab:omc-func}\\
\toprule
\textbf{ID} & \textbf{Trace} & \textbf{Imp.} & \textbf{Units} & \textbf{Marginal} & \textbf{Ideal} \\
\midrule
\endfirsthead
\multicolumn{6}{l}{\small\textit{(Continued from previous page)}}\\[2pt]
\toprule
\textbf{ID} & \textbf{Trace} & \textbf{Imp.} & \textbf{Units} & \textbf{Marginal} & \textbf{Ideal} \\
\midrule
\endhead
\midrule
\multicolumn{6}{r}{\small\textit{(Continued on next page)}}\\
\endfoot
\bottomrule
\endlastfoot
OMC-F-01 & S-01 & C & m$^3$ & $\geq 256$\newline $(8{\times}8{\times}4)$ & Full room volume \\[4pt]
OMC-F-02 & S-01, R-01 & H & m & $\geq 9$ (wide) & $\geq 16$ (narrow) \\[4pt]
OMC-F-03 & S-02, S-04 & H & cameras & $\geq 6$ & $\geq 10$ \\[4pt]
OMC-F-04 & S-01, S-04 & M & m$^2$\,/\,m & $\leq 80$\,m$^2$; ceiling $\leq 5$\,m & Adaptable to room size \\[4pt]
OMC-F-05 & S-04 & M & R & $\leq 50,000$ & Minimised subject to meeting OMC-P-01--06 \\[4pt]
OMC-F-06 & S-04 & L & --- & Mounts removable without specialist tools & Single-person redeployment in $< 4$\,h \\[4pt]
OMC-F-07 & S-04 & L & h & $\geq 2$ & $\geq 4$ \\
\end{longtable}}

%----------------------------------------------------------------------------------------
\section{OCP Subsystem Requirements}
%----------------------------------------------------------------------------------------

The \gls{OCP} subsystem optimises the camera configuration befor physical deployment. Requirements are defined at two levels: \textit{internal} metrics (OCP-P-01--OCP-P-04) characterise solution quality within the simulation model, grounded in the resolution uncertainty model \parencite{sharmaAnalysisUncertaintyBounds1998} and dynamic occlusion model \parencite{rahimianOptimalCameraPlacement2017} of Section~\ref{sec:cost-function}; \textit{external} metrics (OCP-P-05--OCP-P-06) characterise real-world improvement upon deployment, following the validation framework of Section~\ref{sec:verification}.

\subsection{Performance Requirements}

Table~\ref{tab:ocp-perf-desc} defines each metric; quantitative thresholds are given in Table~\ref{tab:ocp-perf}.

{\small\setlength{\tabcolsep}{4pt}
\begin{longtable}{p{1.5cm} >{\raggedright\arraybackslash}p{3.2cm} >{\raggedright\arraybackslash}p{\dimexpr\linewidth-1.5cm-3.2cm-6\tabcolsep\relax}}
\caption[OCP performance metric definitions.]{OCP performance metric definitions. Quantitative thresholds are given in Table~\ref{tab:ocp-perf}.}
\label{tab:ocp-perf-desc}\\
\toprule
\textbf{ID} & \textbf{Metric} & \textbf{Description} \\
\midrule
\endfirsthead
\multicolumn{3}{l}{\small\textit{(Continued from previous page)}}\\[2pt]
\toprule
\textbf{ID} & \textbf{Metric} & \textbf{Description} \\
\midrule
\endhead
\bottomrule
\endlastfoot
OCP-P-01 & Workspace Coverage & Fraction of discrete target-space grid points observed by $\geq 2$ cameras simultaneously. A point visible to fewer than two cameras cannot be triangulated. \\[4pt]
OCP-P-02 & Triangulation Angle Compliance & Fraction of target grid points for which $\geq 1$ camera pair satisfies $40\degree \leq \theta_{ij} \leq 140\degree$. Configurations violating this constraint produce ill-conditioned depth estimates even when the point is visible. \\[4pt]
OCP-P-03 & Combined GA Objective Cost & The weighted combined cost $J = w_{\text{res}}\,J_{\text{res}} + w_{\text{occ}}\,J_{\text{occ}}$ (Eq.~\ref{eq:combined-cost}) that the \gls{GA} directly minimises. Both component costs are normalised so that $J \approx 1$ for a random configuration; values below 1.0 represent improvement over the unoptimised baseline. The two components capture resolution uncertainty (via the pixel-quantisation model) and occlusion robustness (via the dynamic occluder model, Eq.~\ref{eq:occlusion-point}) simultaneously. \\[4pt]
OCP-P-04 & Optimisation Runtime & Time to complete a full run of $G_{\max}$ generations on the available parallel hardware. \\[4pt]
OCP-P-05 & Calibration Residual Improvement & Percentage reduction in the OpenMV system calibration residual (mm) achieved by the optimised configuration relative to the ad-hoc baseline, using identical wand calibration protocol and capture volume. The OptiTrack system serves as an independent ground-truth reference for this comparison. \\[4pt]
OCP-P-06 & Marker Reconstruction Improvement & Increase in marker reconstruction completeness (percentage points) achieved by the optimised vs.\ ad-hoc configuration. Analogous to the 56\,pp improvement demonstrated by Rahimian and Kearney \parencite{rahimianOptimalCameraPlacement2017}. \\
\end{longtable}}

{\small\setlength{\tabcolsep}{4pt}
\begin{longtable}{p{1.5cm} p{2.0cm} p{1cm} p{1.5cm} >{\raggedright\arraybackslash}p{\dimexpr(\linewidth-1.5cm-2.0cm-0.7cm-1.5cm-12\tabcolsep)/2\relax} >{\raggedright\arraybackslash}p{\dimexpr(\linewidth-1.5cm-2.0cm-0.7cm-1.5cm-12\tabcolsep)/2\relax}}
\caption[OCP subsystem performance specifications.]{OCP subsystem performance specifications. Importance: C\,=\,Critical, H\,=\,High, M\,=\,Medium, L\,=\,Low. Metric definitions are given in Table~\ref{tab:ocp-perf-desc}.}
\label{tab:ocp-perf}\\
\toprule
\textbf{ID} & \textbf{Trace} & \textbf{Imp.} & \textbf{Units} & \textbf{Marginal} & \textbf{Ideal} \\
\midrule
\endfirsthead
\multicolumn{6}{l}{\small\textit{(Continued from previous page)}}\\[2pt]
\toprule
\textbf{ID} & \textbf{Trace} & \textbf{Imp.} & \textbf{Units} & \textbf{Marginal} & \textbf{Ideal} \\
\midrule
\endhead
\midrule
\multicolumn{6}{r}{\small\textit{(Continued on next page)}}\\
\endfoot
\bottomrule
\endlastfoot
OCP-P-01 & S-02, S-03, R-04 & C & \% & $> 90$ & $> 99$ \\[4pt]
OCP-P-02 & S-02, S-03, R-05 & \degree & \% & $40 \leq \theta_{ij} \leq 140$ & $60 \leq \theta_{ij} \leq 120$ \\[4pt]
OCP-P-03 & S-03, R-04 & H & (dimensionless) & $< 0.5$ & $< 0.3$ \\[4pt]
OCP-P-04 & S-03, S-04 & H & min & $< 60$ & $< 30$ \\[4pt]
OCP-P-05 & S-02, S-03, R-03 & H & \% & $\geq 20$ & $\geq 50$ \\[4pt]
OCP-P-06 & S-02, S-03, R-04 & H & pp & $\geq 20$ & $\geq 50$ \\
\end{longtable}}

\subsection{Functional Requirements}

Table~\ref{tab:ocp-func-desc} defines each metric; the quantitative thresholds are given in Table~\ref{tab:ocp-func}.

{\small\setlength{\tabcolsep}{4pt}
\begin{longtable}{p{1.5cm} >{\raggedright\arraybackslash}p{3.2cm} >{\raggedright\arraybackslash}p{\dimexpr\linewidth-1.5cm-3.2cm-6\tabcolsep\relax}}
\caption[OCP functional metric definitions.]{OCP functional metric definitions. Quantitative thresholds are given in Table~\ref{tab:ocp-func}.}
\label{tab:ocp-func-desc}\\
\toprule
\textbf{ID} & \textbf{Metric} & \textbf{Description} \\
\midrule
\endfirsthead
\multicolumn{3}{l}{\small\textit{(Continued from previous page)}}\\[2pt]
\toprule
\textbf{ID} & \textbf{Metric} & \textbf{Description} \\
\midrule
\endhead
\bottomrule
\endlastfoot
OCP-F-01 & Solution Repeatability & Consistency of the optimised cost $f^*$ across independent runs with different random seeds. Characterises stochastic variance and bounds the reliability of the solution quality claim. \\[4pt]
OCP-F-02 & Output Interface & Format in which the algorithm exports the optimal configuration $\mathbf{c}^*$ for physical implementation and potential import into the \gls{OMC} software. \\[4pt]
OCP-F-03 & Objective Weight Parameterizability & User-adjustable weighting of $w_{\text{res}}$ and $w_{\text{occ}}$ in the combined cost function (Eq.~\ref{eq:combined-cost}), allowing the designer to bias the optimiser toward precision or occlusion robustness. \\
\end{longtable}}

{\small\setlength{\tabcolsep}{4pt}
\begin{longtable}{p{1.5cm} p{1.8cm} p{0.7cm} p{1.0cm} >{\raggedright\arraybackslash}p{\dimexpr(\linewidth-1.5cm-1.8cm-0.7cm-1.0cm-12\tabcolsep)/2\relax} >{\raggedright\arraybackslash}p{\dimexpr(\linewidth-1.5cm-1.8cm-0.7cm-1.0cm-12\tabcolsep)/2\relax}}
\caption[OCP subsystem functional specifications.]{OCP subsystem functional specifications. Metric definitions are given in Table~\ref{tab:ocp-func-desc}.}
\label{tab:ocp-func}\\
\toprule
\textbf{ID} & \textbf{Trace} & \textbf{Imp.} & \textbf{Units} & \textbf{Marginal} & \textbf{Ideal} \\
\midrule
\endfirsthead
\multicolumn{6}{l}{\small\textit{(Continued from previous page)}}\\[2pt]
\toprule
\textbf{ID} & \textbf{Trace} & \textbf{Imp.} & \textbf{Units} & \textbf{Marginal} & \textbf{Ideal} \\
\midrule
\endhead
\midrule
\multicolumn{6}{r}{\small\textit{(Continued on next page)}}\\
\endfoot
\bottomrule
\endlastfoot
OCP-F-01 & S-03 & M & \%\,CV & $\leq 15$ & $\leq 5$ \\[4pt]
OCP-F-02 & S-03, S-04 & M & --- & Exports $(x,y,z)$ and $(\alpha,\beta,\gamma)$ per camera in a structured file (\texttt{.csv} or \texttt{.mat}) & Directly consumable by the OMC deployment pipeline without manual conversion \\[4pt]
OCP-F-03 & S-03 & M & --- & User-configurable $w_{\text{res}}, w_{\text{occ}} \geq 0$,\newline $w_{\text{res}} + w_{\text{occ}} = 1$ & Extendable to $n$-objective framework \\
\end{longtable}}

%----------------------------------------------------------------------------------------
\section{Constraints and Assumptions}
%----------------------------------------------------------------------------------------

The specifications are subject to the following constraints and assumptions:

\begin{enumerate}

    \item \textbf{Indoor rectangular workspace.} The system is designed for a single indoor room with a flat ceiling and clear floor. The target space is a rectangular volume of approximately $10 \times 9 \times 5$\,m, corresponding to the bounds within the safety net of $[-4, 4] \times [-4, 4] \times [0, 4]$\,m in the world frame.

    \item \textbf{Passive retroreflective marker tracking.} The \gls{OMC} subsystem is designed to track passive infrared-retroreflective markers (15.9\,mm diameter) attached to any platform requiring pose estimation. The system is robot-agnostic; \gls{UAV}s are used as the primary test platform, but the specifications apply to any marker-carrying object within the capture volume. The specifications do not apply to markerless pose estimation or active LED markers.

    \item \textbf{Wall- and ceiling-mounted cameras.} Camera positions are constrained to the four walls and the ceiling (the floor is excluded, since a floor-mounted camera oriented upward contributes negligibly to coverage of the tracking volume and is at risk of being damaged). The mounting bounds match the OCP algorithm search space: $[-5, 5] \times [-4.5, 4.5] \times [0, 4.8]$\,m (Section~\ref{sec:algorithm-config}).

    \item \textbf{Controlled infrared environment.} The specifications assume a laboratory prepared to minimise infrared noise: windows covered or blacked out, highly reflective surfaces minimised, and no competing high-brightness IR sources within the camera fields of view.

    \item \textbf{Static workspace structure.} Obstacles within the capture volume are assumed to be static for the purposes of the \gls{OCP} model, although it is optimised for unknown possible dynamic occlusion from the tracked body itself is captured by the dynamic occluder model component of OCP-P-03, but moving people or additional objects are not modelled.

    \item \textbf{Offline optimisation paradigm.} The \gls{OCP} algorithm is run once prior to physical deployment and does not adapt in real time to camera failures or workspace changes. Re-optimisation is required if the room geometry or camera count changes.

    \item \textbf{Pinhole camera model.} The cost function assumes ideal pinhole projection without lens distortion. In physical deployment, residual post-calibration distortion is assumed negligible relative to the pixel-quantisation uncertainty captured by the resolution model.

\end{enumerate}
 % User requirements + specific goals
\chapter{Optical Motion Capture}
\label{TrackingAlgorithm}

\section{Literature Review}
\subsection{Estimation Algorithms}
\subsection{Hardware Considerations for Camera-Based Sensing}
\subsubsection{Camera Capabilities}
Resolution, latency, iframe rate, microcontroller\dots
\subsubsection{Communication Architectures}
Wired vs wireless: PoE, Wi-Fi, Bluetooth
Synchronous vs Asynchronous

\section{Camera Selection and System Integration}

\section{Theory Development}
\subsection{Rigid-body Motion}
\subsubsection{Reference Frames}
\subsubsection{Representing Pose}

\subsection{Computer Vision}
\subsection{Light and Colour}
\subsubsection{Colour Spaces}
\subsection{Image Formation}
\subsection{Image Processing}
\subsection{Triangulation}

\subsection{Estimation Filtering and Noisy Signals}
\subsubsection{The Kalman Filter}
\subsubsection{The Extended Kalman Filter}
\subsubsection{The Invariant Extended Kalman Filter}
\subsubsection{The Two-Frames Group In-EKF}

\section{Algorithm Development}

\section{Implementation}
\subsection{Camera Calibration}
\subsection{Graphical User Interface}

\section{Experimental Testing Setup}

\section{Results and Discussion}



 % Motion Capture Algorithm + Hardware
\chapter{Optimal Camera Placement}
\label{OCP}

\section{Literature Review}
The quality of a multi-camera localisation system depends on the geometric arrangement of its sensors, and optimising this configuration requires both a model of error sources and an effective search strategy. When the sensor placement is poor, the consequences tend to compound in two distinct ways. The first is measurement uncertainty due to the spatial relationship between camera and target. Cameras discretise their images into pixels; the further a target sits from a camera, the fewer pixels it subtends, and the larger the resulting position estimate. When two cameras observe the same target from nearly parallel viewing angles, the small pixel-level uncertainty propagates into a large displacement along the depth axis, precluding triangulation and consequently the recovered pose estimate. The second degrading factor is visibility loss caused by occlusion (both by obstacles in the environment and the target itself). If a camera is unable to view the target, the amount of measurements are reduced and the uncertainty is amplified \parencite{rahimianOptimalCameraPlacement2017}. Together, these two effects display how an arbitrary camera arrangement can culminate in localisation errors exacerbated 3D recontruction errors beyond the intrinsic noise of the sensors themselves. Thus, optimising the placement is not optional; it should be a prerewuisite for accurate tracking. However, the placement problem is combinatorially explosive, and NP-hard, which necessitates an optimisation strategy capable of exploring a large search space withoin reasonable computational time rather than relying on exhaustive evaluation \parencite{kirchhofOptimalPlacementMultiple2013}. This section reviews the literature on the \gls{OCP} problem for both of these factes: first, the optimisation strategies employed to search the multi-objective, unsmooth, ill-defined search space in sufficently good time, from heuristic algorithms to integer programming formulations; and second, how the sources of meausrement uncertainty, quantisation and occlusion, have been characterised and modelled.  


\subsection{Background of OCP}
The OCP descends from the well-studied \gls{AGP}, which investigated the minimum number of guards required to observe an entire 2D polygon with $n$ vertices. The AGP assumed guards had unlimited $360 \degree$ field of view, infinite range, and operated in two dimensions only. Fisk elegantly proved that $\lfloor\frac{n}{3}\rfloor$ guards are always sufficient to cover any simple polygon, using a proof based on triangulation and 3-colouring; in practice, specific polygons often require fewer, as demonstrated in Figure %\ref{fig:fisk} \parencite{fiskArtGalleryProblem1978}. 
In this example, a polygon with $n=13$ vertices is first triangulated, then each vertex is assigned one of three colours such that no triangle shares two vertices of the same colour; placing guards at all vertices of the least frequent colour guarantees complete coverage. Similar to the AGP the OCP must take a heauristic goal and mathematically model it in pursuit that given the amount of sensors the  arrangment. The optimal camera placement problem extends this foundation by introducing constraints that reflect the true capabilities and limitations of physical sensors. Moreover, the OCP seeks to extend beyond mere coverage as in the case of accurate triangulation more than one camera needs to view a point in the target area. 

\subsection{Motivation for Optimal Placement}
These constraints are motivated by diverse applications, such as: surveillance and monitoring \parencite{bodorOptimalCameraPlacement2007}, motion capture \parencite{aissaouiDesigningCameraPlacement2018, zieglerOptimalCameraPlacement2023}, augmented and virtual reality \parencite{rahimianOptimalCameraPlacement2017}, and autonomous driving \parencite{puligandlaMultiresolutionApproachLarge2022}. In such visual sensor networks, accurate localisation requires that each marker or feature be visible to at least two cameras, enabling 3D position recovery through triangulation \parencite{aissaouiDesigningCameraPlacement2018,zieglerOptimalCameraPlacement2023}. Most papers emphasise that sensor placement (including amount, location, orientation, and arrangement to the other sensors in the network) significantly affects system coverage and tracking performance \parencite{zhouSensorPlacementOptimization2024}. However, the OCP is NP-hard and combinatorial; this means that the search space grows explosively with the number of candidate positions and orientations, rendering exhaustive search intractable. Consequently, most investigators resort to heuristic-based approximation algorithms as they offer a practical path to reach acceptable solutions within reasonable computation time.

Generally, the OCP can be stated as follows: given a set of candidate camera locations and orientations, find the configuration that maximises 3D reconstruction quality of a target within a defined workspace, subject to practical constraints such as mounting options and budget. Solving this problem requires balancing trade-offs among coverage, accuracy, and computational efficiency. Without the aid of an optimisation algorithm, the sensors are arranged in an ad-hoc configuration by the designer, which requires expertise or trial and error, which can be time-consuming or sacrificial if a sub-par arrangement is settled on with no mathematical  quantitative guidance. Thus, the resulting localisation error is neither calculated nor considered mathematically. Here lies the need to optimise the camera network based on metrics that minimise degrading effects like resolution uncertainty and probabilistic occlusion even if the trajectory of the target is not known exactly beforehand.

\subsection{Optimisation Approaches for Sensor Placement}
Recent research has diverged into avenues of how to formulate the problem, when to conclude that an optimal solution has been found, and what consitutes success as a benchmark. Two primary classes of optimisation methods have emerged: heuristic algorithms and integer programming. The problem formulation typically depends of defining a specific objective function, sensor model, and workspace representation. Eventhough the search space is continuous it is typical of designers to discretise both the workspace and the possible camera locations to make the problem more computable. Depending on the application, available resources, and workspace constraints, different optimisation objectives have been prioritised in the literature. Broadly, these can be categorised into coverage-based objectives and accuracy-based objectives. In terms of coverage-based objectives, the goal is to maximise the visible area of the workspace, and sometimes going further to handle aspects of occlusion both static and dynamic and hence usually ensuring a two-cover coverage. This is particularly relevant in surveillance applications where ensuring that all regions are monitored is critical \parencite{bodorOptimalCameraPlacement2007}. Conversely, accuracy-based objectives focus on minimising localisation error, which is crucial in motion capture and augmented reality applications where precise tracking is required and more than one camera arranged sufficiently non-parallel is needed to ensure triangulation is possible \parencite{aissaouiDesigningCameraPlacement2018, rahimianOptimalCameraPlacement2017}. Wu, Sharma, and Huang proposed another accuracy-based metric of placement which assessed how the limited camera intrinsics and quantisation lead to inaccurate measurements exacerbated by the placement of two cameras \parencite{sharmaAnalysisUncertaintyBounds1998}. Many studies have proposed hybrid objectives that balance the trade-off of accuracy with coverage \parencite{zieglerOptimalCameraPlacement2023, zhouSensorPlacementOptimization2024}. Some works have also considered additional constraints such as budget limits and physical mounting options \parencite{aissaouiDesigningCameraPlacement2018, puligandlaMultiresolutionApproachLarge2022, daiOptimalCameraConfiguration}.

\subsubsection{Optimisation Objectives}
The configuration quality for 3D marker tracking and trinagulation is mostly affected by two main sources of error: marker visibility and triangulation accuracy. In terms of the application being \gls{MoCap}, a marker should not be occluded by any obstacle, including itself, the marker must be within the cameras field of view, vissible given the camera's resolution and range of capture. These are necessary to ensure consistent and continual tracking. Furthermore, at least two cameras should be arranged sufficiently non-parallel so that triangulation is well-conditioned. Sometimes, the easy answer is just to throw more cameras at the problem as that does give redundant views and ensures we can place camera pairs to cover the scene in a sufficently non-parallel format but in most research situationships, it is not prudent to over-exhaust the budget and hence it may be a desirable goal to asses the least amount of cameras required to achieve a level of accuracy/ coverage to minimise cost.  

Provided below is a summary table of the different optimisation qualities found across papers in the literature
(insert table with rows indication paper citation and application and columns of optimsation objectives)

\subsection{Modelling the Problem}
As mentioned earlier due to the large search space which is made so large becaus ethe problem is combinatorial (meaning additonal combinations of different aspects like position and roientation of each camera) which is NP-hard and increases expeonently with each additonal sensor added to the network it becomes necessary to make some simplifications as an exhaustive search especially for a room-sized scale search simply becomes intractable and computationally inefficent. This is often why the search-space is dicsretised both in the flight-space arena and in the possible mounting locations of the camera. 

\subsection{Summary}

The measurement uncertainty attributed to the spatial relationship between the target and the cameras can be modelled based on the distance between the target and camera which can excerbated pixel resolution and the orientation between camera's view vectors where the increasingly more parallel these vectors are results in ill-conditioned triangulation geometry, culminating in poor pose estimates and hence poor tracking

As seen there have been a plethora of proposed methodologies for optimising camera placement in the literature which are summarised in the timeline below where they have been broadly categorised as heuristic optimisation or binary integer programming and their objectives have also been listed based on thematic colour rings indicated in the key. 

\section{Theoretical Framework}
\subsection{Genetic Algorithm Terminology}
The language used to describe genetic algorithms borrows heavily from evolutionary biology. Table~\ref{tab:ga-terminology} defines these terms in both domains to clarify the analogy and establish consistent vocabulary for the sections that follow.

\begin{table}[h]
\centering
\caption{Comparison of biological and genetic algorithm terminology.}
\begin{tabularx}{\textwidth}{l X X}
\toprule 
Term & Biological Meaning & Genetic Algorithm Meaning \\
\midrule 
Chromosome & A structure of \gls{DNA} carrying genetic information & A complete candidate solution encoded in some representation, e.g bit strings \\
Gene & A unit of heredity determining a specific trait & A single element or parameter within the candidate solution, e.g a bit \\
Population & A group of interbreeding individuals of the same species & The set of candidate solutions existing at a given iteration \\
Fitness & An organism's ability to survive and reproduce in its environment & An objective function quantifying how well a candidate solution achieves the desired goal \\
Generation & All individuals born and living at approximately the same time & One complete iteration of the algorithm producing a new population \\
Selection & Differential survival and reproduction due to heritable trait variation & An operator that chooses higher-performing individuals to serve as parents \\
Crossover & Exchange of genetic material between homologous chromosomes during meiosis & An operator that combines portions of two parent solutions to create offspring \\
Mutation & A random, heritable change in genetic sequence, ususally caused by a erroneous \gls{DNA} replication & An operator that randomly perturbs one or more genes in a solution \\
\bottomrule
\end{tabularx}
\label{tab:ga-terminology}
\end{table}

This biological analogy provides useful intuition: just as natural selection refines a species over generations by favouring advantageous traits, genetic algorithms iteratively improve a population of solutions by favouring those with higher fitness. While genetic algorithms lack the biological complexity of true evolution, they do provide a productive conceptual framework. The effectiveness of this process depends on how solutions are represented, this is often reffered to as the encoding.
\subsection{Encoding Schemes}
\subsection{Selection Methods}
\subsection{Crossover Techniques}
\subsection{Convergence Analysis}



\section{Methodology}
\subsection{Algorithm Configuration}
\subsubsection{Parameter Selection}
\subsubsection{Initial Population}
\subsubsection{Tournament Selection}
\subsubsection{Double Point Crossover}
\subsubsection{Mutation}
\subsubsection{Elitism}
\subsection{Cost Function}
\subsubsection{Resolution Uncertainty Cost Function}
\subsubsection{Dynamic Occlusion Cost Function}
\subsubsection{Combined Cost Function}
\subsection{Termination Criteria}


\section{Results and Discussion}
\subsection{Convergence Behaviour}
\subsection{Cost Function Evaluation}
\subsection{Parameter Sensitivity Analysis}
\subsection{Optimal Placement Configurations}
Different Modalities
Random Placement, ad hoc placement 
\subsection{System Limitations} ? % Optimal Camera Placement and Genetic Algorithm 
\chapter{Conclusion and Future Work}
\label{Conclusion}
 % Conclusions



\chapterstyle{default} % Reset the chapter style back to the default used for non-content chapters

%----------------------------------------------------------------------------------------
%	REFERENCES
%----------------------------------------------------------------------------------------
\printbibliography[title={References}] % Uses biblatex/biber

%----------------------------------------------------------------------------------------
% APPENDICES
%----------------------------------------------------------------------------------------
\chapterstyle{thesis} % Change the style of the Chapter header to that defined in structure.tex
\include{Chapters/Appendices}

%----------------------------------------------------------------------------------------
%	INDEX
%----------------------------------------------------------------------------------------

\printindex % Print the index

%----------------------------------------------------------------------------------------

\end{document}